% ============================================================
%  C: Como Programar -- 6ª Edição | Deitel & Deitel
%  Capítulo 1 -- Introdução aos Computadores, à Internet e à Web
%  FAZENDO A DIFERENÇA (1.10, 1.11, 1.12)
%  Pesquisa externa + raciocínio algorítmico introdutório
%  Estilo didático e incremental -- no espírito de Paul Deitel
% ============================================================
\documentclass[12pt,a4paper]{article}

\usepackage[utf8]{inputenc}
\usepackage[T1]{fontenc}
\usepackage[brazil]{babel}
\usepackage{geometry}
\geometry{top=2.5cm, bottom=2.5cm, left=3cm, right=3cm}
\usepackage{parskip}
\setlength{\parskip}{6pt}
\setlength{\parindent}{0pt}


\usepackage{xcolor}
\usepackage{tcolorbox}
\tcbuselibrary{skins, breakable, theorems}
\usepackage{enumitem}
\usepackage{listings}
\usepackage{fancyhdr}
\usepackage{titlesec}
\usepackage{hyperref}
\usepackage{booktabs}
\usepackage{array}
\usepackage{amsmath}

% Paleta Deitel
\definecolor{deitelblue}{RGB}{0, 70, 127}
\definecolor{deitelgray}{RGB}{240, 242, 245}
\definecolor{deitelgreen}{RGB}{0, 110, 60}
\definecolor{deitellorange}{RGB}{200, 90, 0}
\definecolor{deitellightblue}{RGB}{220, 235, 250}
\definecolor{deitelred}{RGB}{160, 20, 20}
\definecolor{ecogreen}{RGB}{0, 120, 50}
\definecolor{ecolightgreen}{RGB}{210, 240, 220}
\definecolor{saude}{RGB}{150, 0, 60}
\definecolor{saudelightred}{RGB}{245, 215, 225}
\definecolor{genero}{RGB}{100, 0, 150}
\definecolor{generolightpurple}{RGB}{235, 220, 245}

% Caixas
\newtcolorbox{boaprática}[1][]{colback=deitellightblue,colframe=deitelblue,
  fonttitle=\bfseries\small,title={Boa prática de programação},breakable,#1}
\newtcolorbox{pesquisa}[1][]{colback=ecolightgreen,colframe=ecogreen,
  fonttitle=\bfseries\small,title={Pesquisa externa realizada},breakable,#1}
\newtcolorbox{respostabox}[1][]{colback=deitelgray,colframe=deitelblue!60,
  fonttitle=\bfseries\small\color{deitelblue},title={Resposta detalhada},breakable,#1}
\newtcolorbox{enunciadoBox}[1][]{colback=white,colframe=deitelblue,
  fonttitle=\bfseries,title={#1},breakable}
\newtcolorbox{algoritmo}[1][]{colback=deitelgray,colframe=deitellorange,
  fonttitle=\bfseries\small\color{deitellorange},title={Descrição do Algoritmo},breakable,#1}
\newtcolorbox{atencao}[1][]{colback=yellow!10,colframe=deitellorange,
  fonttitle=\bfseries\small,title={Atenção},breakable,#1}
\newtcolorbox{ecoCaixa}[1][]{colback=ecolightgreen,colframe=ecogreen,
  fonttitle=\bfseries,title={#1},breakable}
\newtcolorbox{saudeCaixa}[1][]{colback=saudelightred,colframe=saude,
  fonttitle=\bfseries,title={#1},breakable}
\newtcolorbox{generoCaixa}[1][]{colback=generolightpurple,colframe=genero,
  fonttitle=\bfseries,title={#1},breakable}

\setlist[enumerate,1]{label=\textbf{\alph*)},leftmargin=2em}
\setlist[itemize]{leftmargin=2em}

\lstset{
  basicstyle=\ttfamily\small,
  backgroundcolor=\color{deitelgray},
  frame=single,framerule=0.4pt,
  rulecolor=\color{deitelblue!40},
  breaklines=true,showstringspaces=false,
  keywordstyle=\color{deitelblue}\bfseries,
  commentstyle=\color{ecogreen}\itshape,
  stringstyle=\color{deitelred},
  language=C,numbers=none,
  xleftmargin=10pt,xrightmargin=5pt,
}

\pagestyle{fancy}
\fancyhf{}
\fancyhead[L]{\small\color{deitelblue}\textbf{C: Como Programar -- 6ª ed.\ $|$ Deitel \& Deitel}}
\fancyhead[R]{\small\color{deitelblue}Cap.\ 1 -- Fazendo a Diferença}
\fancyfoot[C]{\small\thepage}
\renewcommand{\headrulewidth}{0.4pt}

\titleformat{\section}[block]{\large\bfseries\color{deitelblue}}{\thesection.}{0.5em}{}[\titlerule]
\titleformat{\subsection}[block]{\normalsize\bfseries\color{deitelblue!80}}{\thesubsection.}{0.5em}{}

\hypersetup{colorlinks=true,linkcolor=deitelblue,urlcolor=ecogreen}

% ============================================================
\begin{document}

% ---- Capa ----
\begin{titlepage}
  \centering
  \vspace*{2cm}
  {\color{deitelblue}\rule{\linewidth}{2pt}}\\[0.4cm]
  {\huge\bfseries\color{deitelblue} C: Como Programar}\\[0.2cm]
  {\large\color{deitelblue} 6ª Edição $\cdot$ Paul Deitel \& Harvey Deitel}\\[0.2cm]
  {\color{deitelblue}\rule{\linewidth}{2pt}}\\[1cm]
  {\LARGE\bfseries Capítulo 1}\\[0.3cm]
  {\Large\color{deitelblue!80} Introdução aos Computadores, à Internet e à Web}\\[0.8cm]
  \begin{tcolorbox}[colback=ecolightgreen,colframe=ecogreen,width=0.85\linewidth]
    \centering
    {\large\bfseries\color{ecogreen} Fazendo a Diferença}\\[0.3cm]
    {\normalsize Exercícios 1.10 $\cdot$ 1.11 $\cdot$ 1.12}\\[0.2cm]
    {\small Inclui pesquisa externa, análise de ferramentas e\\
    raciocínio algorítmico introdutório}
  \end{tcolorbox}
  \vfill
  {\small\color{deitelblue} ``Encorajamos você a usar computadores e a Internet para pesquisar e
  resolver\\problemas que realmente fazem a diferença.'' -- Deitel \& Deitel, Prefácio}
\end{titlepage}

% ---- Introdução ----
\section*{Sobre os Exercícios ``Fazendo a Diferença''}

A série \textbf{Fazendo a Diferença} foi introduzida nesta 6ª edição como um conjunto
especial de exercícios que vai além do domínio técnico de C. Nosso objetivo é
encorajá-lo a usar computadores e a Internet para pesquisar e resolver problemas que
realmente importam para a sociedade e para o mundo. Esperamos que você encare esses
exercícios com seus próprios valores, crenças e visão de mundo.

Diferentemente dos exercícios anteriores, alguns itens desta seção envolvem
\textbf{pesquisa na Web} e uma abordagem mais reflexiva e descritiva -- afinal,
no Capítulo~1 ainda não aprendemos a programar em C. O que podemos fazer, com toda
a profundidade, é \textbf{descrever algoritmos} -- isto é, os passos lógicos e
ordenados que um futuro programa seguiria para resolver cada problema.

\begin{atencao}
  Lembre-se: no Capítulo~1, não escrevemos código C. Descrevemos os problemas,
  exploramos as ferramentas disponíveis e formulamos o raciocínio algorítmico
  que será implementado em capítulos posteriores. No Capítulo~3 você aprenderá
  o que é um algoritmo formalmente, e no Capítulo~4 aprenderá a expressá-lo
  em pseudocódigo e em C.
\end{atencao}

\newpage
\tableofcontents
\newpage

% ============================================================
\section{Exercício 1.10 -- \textit{Test-Drive}: Calculadora de Emissão de Carbono}
% ============================================================

\begin{enunciadoBox}[Enunciado]
\textbf{\textit{Test-drive}: Calculadora de emissão de carbono.}
Alguns cientistas acreditam que as emissões de carbono, especialmente as ocorridas pela
queima de combustíveis fósseis, contribuem de modo significativo para o aquecimento global,
e que isso pode ser combatido se os indivíduos tomarem medidas que limitem seu uso de
combustíveis com base no carbono. Sites da Web como o TerraPass
(\texttt{www.terrapass.com/carbon-footprint-calculator/}) e o Carbon Footprint
(\texttt{www.carbonfootprint.com/calculator.aspx}) oferecem calculadoras de emissão de
carbono. \textbf{Faça o teste dessas calculadoras para estimar sua emissão de carbono.}
Para se preparar, pesquise na Web as fórmulas usadas no cálculo.
\end{enunciadoBox}

\bigskip

\begin{pesquisa}
Este exercício requer pesquisa externa, conforme solicitado pelo autor. As informações
a seguir foram obtidas consultando as calculadoras indicadas e a literatura científica
sobre pegada de carbono (\textit{carbon footprint}).
\end{pesquisa}

\bigskip

\subsection*{1.1 \ O que é Pegada de Carbono?}

A \textbf{pegada de carbono} (\textit{carbon footprint}) de um indivíduo ou organização
é a quantidade total de \textbf{gases de efeito estufa (GEE)} -- principalmente dióxido
de carbono (CO\textsubscript{2}) e metano (CH\textsubscript{4}) -- emitidos como resultado
direto ou indireto de suas atividades. É medida em \textbf{toneladas de CO\textsubscript{2}
equivalente} (tCO\textsubscript{2}e) por ano.

\bigskip

\subsection*{1.2 \ Principais Fontes de Emissão Individual}

\begin{center}
\renewcommand{\arraystretch}{1.5}
\begin{tabular}{l l p{6.5cm}}
  \toprule
  \textbf{Categoria} & \textbf{Fonte} & \textbf{Fator de Emissão (referência)}\\ \midrule
  Transporte & Automóvel a gasolina &
    $\approx 2{,}31$ kg CO\textsubscript{2}/litro de gasolina\\
  Transporte & Voos domésticos &
    $\approx 0{,}255$ kg CO\textsubscript{2}/km/passageiro\\
  Energia & Eletricidade (Brasil -- fator médio 2023) &
    $\approx 0{,}100$ kg CO\textsubscript{2}/kWh (matriz predominantemente hidrelétrica)\\
  Alimentação & Carne bovina &
    $\approx 27$ kg CO\textsubscript{2}e/kg consumido\\
  Alimentação & Dieta vegetariana &
    $\approx 1{,}5$ kg CO\textsubscript{2}e/dia\\
  Consumo & Compras gerais &
    Varia amplamente; estimativa: 0,5--2,0 tCO\textsubscript{2}e/ano\\
  \bottomrule
\end{tabular}
\end{center}

\bigskip

\subsection*{1.3 \ Análise das Calculadoras Testadas}

\textbf{a) TerraPass Carbon Footprint Calculator}
(\url{https://www.terrapass.com/carbon-footprint-calculator/})

A calculadora da TerraPass divide a análise em categorias:
\begin{itemize}
  \item \textbf{Veículo:} solicita milhas percorridas por ano e o consumo do veículo
  (milhas por galão -- MPG). A fórmula subjacente é:
  \[
    \text{CO}_2\ (\text{libras}) =
    \frac{\text{milhas anuais}}{\text{MPG}} \times 19{,}6
  \]
  onde o fator $19{,}6$ é o número de libras de CO\textsubscript{2} emitidas por galão
  de gasolina queimada (EPA -- Agência de Proteção Ambiental dos EUA).

  \item \textbf{Voos:} solicita número de voos de curta, média e longa distância.
  Cada trecho tem um fator de emissão estimado em kg CO\textsubscript{2} por passageiro.

  \item \textbf{Residência:} consumo de gás natural, óleo combustível e eletricidade.
\end{itemize}

\textbf{b) Carbon Footprint Calculator}
(\url{https://www.carbonfootprint.com/calculator.aspx})

Esta calculadora é mais abrangente e inclui:
\begin{itemize}
  \item Eletricidade e gás doméstico
  \item Voos (domésticos e internacionais, com classes diferentes de cabine)
  \item Transporte público, carro pessoal e motocicleta
  \item Estilo de vida: alimentação, compras, serviços, resíduos
\end{itemize}

A metodologia segue o protocolo GHG (\textit{Greenhouse Gas Protocol}) da WBCSD/WRI,
padrão internacional para contabilidade de emissões.

\bigskip

\subsection*{1.4 \ Simulação de Teste: Perfil Típico de Estudante Universitário Brasileiro}

Para ilustrar o uso da calculadora, realizamos uma estimativa com o seguinte perfil:

\begin{center}
\renewcommand{\arraystretch}{1.4}
\begin{tabular}{l r}
  \toprule
  \textbf{Atividade} & \textbf{Emissão estimada (tCO\textsubscript{2}e/ano)} \\ \midrule
  Transporte de carro (15.000 km/ano, carro flex) & $\approx 1{,}50$ \\
  Eletricidade residencial (200 kWh/mês) & $\approx 0{,}24$ \\
  Alimentação (onívoro moderado) & $\approx 1{,}20$ \\
  Sem voos internacionais & $0{,}00$ \\
  Consumo geral & $\approx 0{,}80$ \\ \midrule
  \textbf{Total aproximado} & $\approx \mathbf{3{,}74}$ \\ \bottomrule
\end{tabular}
\end{center}

\medskip
Para comparação, a média mundial é de $\approx 4{,}7$ tCO\textsubscript{2}e/pessoa/ano,
e o limite sustentável para 1,5°C de aquecimento é de aproximadamente
$\mathbf{2{,}3}$ tCO\textsubscript{2}e/pessoa/ano (IPCC, 2018).

\bigskip

\subsection*{1.5 \ Esboço do Algoritmo (preparação para o Capítulo~2)}

Ainda que no Capítulo~1 não escrevamos código C, podemos esboçar o \textbf{algoritmo}
(sequência lógica de passos) que um programa de calculadora de emissões seguiria.
Você implementará isso em capítulos posteriores:

\begin{algoritmo}
\textbf{Algoritmo: Calculadora de Emissão de Carbono}

\begin{enumerate}
  \item Obter do usuário: quilômetros percorridos de carro por ano e consumo médio
  do veículo (km/litro).
  \item Calcular emissão do carro:
  \[
    E_{\text{carro}} = \frac{\text{km}}{\text{km/litro}} \times 2{,}31
    \quad [\text{kg CO}_2]
  \]
  \item Obter do usuário: consumo mensal de eletricidade (kWh).
  \item Calcular emissão de eletricidade:
  \[
    E_{\text{eletricidade}} = \text{kWh\_mensal} \times 12 \times \text{fator\_regional}
    \quad [\text{kg CO}_2]
  \]
  \item Obter do usuário: tipo de dieta (carnívora, onívora, vegetariana, vegana).
  \item Associar dieta a um fator de emissão anual pré-determinado.
  \item Somar todas as emissões e converter para toneladas:
  \[
    E_{\text{total}} = (E_{\text{carro}} + E_{\text{eletricidade}} + E_{\text{dieta}} + \ldots)
    \div 1000 \quad [\text{tCO}_2\text{e/ano}]
  \]
  \item Exibir o resultado e compará-lo com médias nacionais e globais.
  \item Sugerir ações para reduzir a pegada de carbono.
\end{enumerate}
\end{algoritmo}

\begin{boaprática}
  Exercícios futuros (Capítulos~2 e~4) pedirão que você implemente em C uma
  calculadora de emissão de carbono. Este exercício de \textit{test-drive} o
  prepara conceitualmente: ao usar as calculadoras existentes, você entende
  quais dados são necessários, quais cálculos são feitos e qual seria a saída
  esperada do seu programa. Essa compreensão do problema \textit{antes} de
  codificá-lo é uma excelente prática de engenharia de software.
\end{boaprática}

\newpage

% ============================================================
\section{Exercício 1.11 -- \textit{Test-Drive}: Calculadora de IMC}
% ============================================================

\begin{enunciadoBox}[Enunciado]
\textbf{\textit{Test-drive}: Calculadora de índice de massa corporal.}
De acordo com estimativas recentes, dois terços da população dos EUA estão acima do
peso, e cerca de metade dela é obesa. Isso causa aumentos significativos na incidência
de doenças como diabetes e problemas cardíacos. Para determinar se uma pessoa está com
excesso de peso ou sofre de obesidade, pode-se usar o \textbf{Índice de Massa Corporal
(IMC)}. Nos EUA, o Department of Health and Human Services oferece uma calculadora
de IMC em \texttt{www.nhlbisupport.com/bmi/}. Use-a para calcular seu IMC. Um
exercício no Capítulo~2 lhe pedirá para programar sua própria calculadora de IMC.
Para se preparar, pesquise na Web as fórmulas utilizadas no cálculo.
\end{enunciadoBox}

\bigskip

\begin{pesquisa}
Este exercício requer pesquisa externa. As informações a seguir foram obtidas com
base nos critérios do WHO (Organização Mundial da Saúde), NIH (Institutos Nacionais
de Saúde dos EUA) e Ministério da Saúde do Brasil.
\end{pesquisa}

\bigskip

\subsection*{2.1 \ O que é o Índice de Massa Corporal (IMC)?}

O \textbf{IMC} (\textit{Body Mass Index} -- BMI) é uma medida internacional
amplamente utilizada para avaliar se o peso de uma pessoa está dentro de uma
faixa saudável, considerando sua altura. É uma medida de triagem (não de diagnóstico),
simples de calcular e aplicável a adultos.

\bigskip

\subsection*{2.2 \ Fórmula do IMC}

\textbf{Sistema métrico (padrão utilizado no Brasil):}
\[
  \boxed{IMC = \frac{\text{massa} \, [\text{kg}]}{\text{altura}^2 \, [\text{m}^2]}}
\]

\textbf{Exemplo:} Uma pessoa com 75 kg e 1,75 m de altura:
\[
  IMC = \frac{75}{1{,}75^2} = \frac{75}{3{,}0625} \approx 24{,}5 \, \text{kg/m}^2
\]

\textbf{Sistema imperial (usado nas calculadoras americanas):}
\[
  IMC = \frac{\text{peso} \, [\text{libras}] \times 703}{\text{altura}^2 \, [\text{polegadas}^2]}
\]

O fator 703 é o fator de conversão de unidades imperiais para o sistema métrico.

\bigskip

\subsection*{2.3 \ Classificação da OMS por Faixa de IMC}

\begin{center}
\renewcommand{\arraystretch}{1.5}
\begin{tabular}{c l l}
  \toprule
  \textbf{IMC (kg/m²)} & \textbf{Classificação (OMS)} & \textbf{Risco associado} \\ \midrule
  $< 18{,}5$ & Abaixo do peso & Baixo (outros riscos) \\
  $18{,}5$ -- $24{,}9$ & \textcolor{ecogreen}{\bfseries Peso normal} & Médio \\
  $25{,}0$ -- $29{,}9$ & Sobrepeso & Aumentado \\
  $30{,}0$ -- $34{,}9$ & Obesidade Grau I & Alto \\
  $35{,}0$ -- $39{,}9$ & Obesidade Grau II & Muito alto \\
  $\geq 40{,}0$ & Obesidade Grau III (mórbida) & Extremamente alto \\ \bottomrule
\end{tabular}
\end{center}

\medskip
\textbf{Nota importante:} Para crianças e adolescentes (menores de 18 anos), o IMC
é interpretado por meio de percentis por idade e sexo, não pelas faixas acima.
Para idosos, os limites podem ser ligeiramente diferentes.

\bigskip

\subsection*{2.4 \ Limitações do IMC}

O IMC é uma medida útil de triagem, porém possui limitações conhecidas:
\begin{itemize}
  \item \textbf{Não distingue massa muscular de gordura:} Atletas com alta massa
  muscular podem ter IMC elevado sem excesso de gordura.
  \item \textbf{Não indica distribuição de gordura:} A gordura abdominal (visceral)
  é mais perigosa que a gordura subcutânea, mas o IMC não as diferencia.
  \item \textbf{Variações étnicas:} Algumas populações (asiáticas, por exemplo)
  apresentam maior risco cardiometabólico com IMC mais baixo.
\end{itemize}

\bigskip

\subsection*{2.5 \ Algoritmo para a Calculadora de IMC (preparação para o Capítulo~2)}

O exercício menciona que o Capítulo~2 pedirá que você \textit{programe} sua calculadora
de IMC. Veja o algoritmo que o futuro programa C deverá seguir:

\begin{algoritmo}
\textbf{Algoritmo: Calculadora de IMC}

\begin{enumerate}
  \item Exibir mensagem pedindo o peso do usuário em quilogramas.
  \item Ler o valor do peso (número real positivo).
  \item Exibir mensagem pedindo a altura do usuário em metros.
  \item Ler o valor da altura (número real positivo).
  \item Calcular o IMC:
  \[
    imc \leftarrow peso \div (altura \times altura)
  \]
  \item Exibir o valor calculado do IMC.
  \item Comparar o IMC com os limites de classificação e exibir a categoria:
  \begin{itemize}
    \item Se $imc < 18{,}5$ $\rightarrow$ exibir ``Abaixo do peso''
    \item Se $18{,}5 \leq imc < 25{,}0$ $\rightarrow$ exibir ``Peso normal''
    \item Se $25{,}0 \leq imc < 30{,}0$ $\rightarrow$ exibir ``Sobrepeso''
    \item Se $imc \geq 30{,}0$ $\rightarrow$ exibir ``Obesidade''
  \end{itemize}
  \item Exibir recomendação de consultar um profissional de saúde.
\end{enumerate}
\end{algoritmo}

\begin{boaprática}
  Este algoritmo será implementado em C no Capítulo~2, usando os comandos
  \texttt{printf} (para exibir mensagens), \texttt{scanf} (para ler dados) e
  a estrutura de decisão \texttt{if/else} (para classificar o IMC), que você
  aprenderá no Capítulo~3. Ao entender o problema completamente \textit{antes}
  de programar, você estará seguindo a melhor prática de engenharia de software:
  \textit{``Pense antes de codificar''}.
\end{boaprática}

\subsection*{2.6 \ Exemplo de Saída Esperada do Futuro Programa}

Quando implementado nos capítulos seguintes, o programa produzirá uma saída similar a:

\begin{tcolorbox}[colback=black!90, colframe=deitelblue, fontupper=\ttfamily\small\color{green!70}]
Digite seu peso em kg: \textcolor{white}{75.0}\\
Digite sua altura em metros: \textcolor{white}{1.75}\\
\\
Seu IMC e: 24.49\\
Classificacao: Peso normal\\
Parabens! Seu peso esta dentro da faixa saudavel recomendada pela OMS.\\
Consulte sempre um profissional de saude para orientacoes personalizadas.
\end{tcolorbox}

\newpage

% ============================================================
\section{Exercício 1.12 -- Neutralidade de Gênero na Linguagem}
% ============================================================

\begin{enunciadoBox}[Enunciado]
\textbf{Neutralidade de gênero.} Muitas pessoas desejam eliminar o sexismo em todas
as formas de comunicação. Foi pedido que você crie um programa que possa processar
um parágrafo de texto e substituir palavras com gênero específico por outras sem
gênero. Supondo que você tenha recebido uma lista de palavras de gênero específico
e suas substitutas com gênero neutro (por exemplo, ``esposa'' por ``cônjuge'',
``homem'' por ``pessoa'', ``filha'' por ``prole'', e assim por diante), explique
o \textbf{procedimento} que você usaria para ler um parágrafo de texto e realizar
essas substituições manualmente. Como seu procedimento trataria um termo como
``super-homem''? Você descobrirá no Capítulo~4 que o termo mais formal para
``procedimento'' é \textbf{algoritmo}, e que um algoritmo especifica as etapas a
serem cumpridas e a ordem em que se deve realizá-las.
\end{enunciadoBox}

\bigskip

\subsection*{3.1 \ Contexto e Motivação}

Este exercício aborda um problema de \textbf{processamento de texto} -- uma das
aplicações mais comuns da programação. Embora o Capítulo~1 não nos ensine a
escrever código C, ele nos pede para pensar algoritmicamente: como um ser humano
\textit{(ou futuramente um programa)} leria um texto e realizaria substituições
de palavras sensíveis ao gênero?

Este tipo de problema é relevante em:
\begin{itemize}
  \item Processadores de texto e corretores automáticos
  \item Ferramentas de revisão de documentos corporativos e acadêmicos
  \item Sistemas de tradução e localização de software
\end{itemize}

\bigskip

\subsection*{3.2 \ O Dicionário de Substituições}

O primeiro passo é ter um \textbf{dicionário de substituições} -- uma lista de pares
\textit{(palavra\_original $\to$ palavra\_neutra)}. Exemplo:

\begin{center}
\renewcommand{\arraystretch}{1.4}
\begin{tabular}{l l}
  \toprule
  \textbf{Palavra com gênero} & \textbf{Substituta neutra} \\ \midrule
  esposa / marido & cônjuge \\
  homem / mulher & pessoa \\
  filho / filha & filho(a) / prole \\
  pai / mãe & progenitor(a) / responsável \\
  ator / atriz & artista \\
  presidente / presidenta & presidente \\
  homens & pessoas \\
  super-homem & super-herói / pessoa extraordinária \\
  \bottomrule
\end{tabular}
\end{center}

\bigskip

\subsection*{3.3 \ O Procedimento (Algoritmo) de Substituição Manual}

Veja o algoritmo que uma pessoa (ou futuramente um programa C) seguiria:

\begin{algoritmo}
\textbf{Algoritmo: Substituição de Palavras com Gênero -- Versão Simples}

\begin{enumerate}
  \item \textbf{Preparar o dicionário:} Ter em mãos (ou na memória do computador)
  a lista completa de pares \textit{(palavra\_original $\to$ palavra\_neutra)}.

  \item \textbf{Ler o texto:} Ler o parágrafo de texto do início ao fim,
  palavra por palavra (ou caractere por caractere, no caso de um computador).

  \item \textbf{Para cada palavra do texto, repetir:}
  \begin{enumerate}[label=\arabic*.]
    \item Isolar a palavra atual (identificar onde ela começa e termina --
    delimitada por espaços, pontuação, quebras de linha).
    \item \textit{Remover} a pontuação temporariamente (vírgulas, pontos, etc.)
    para comparar apenas as letras.
    \item Comparar a palavra (em letras minúsculas, para ignorar maiúsculas)
    com cada entrada no dicionário de substituições.
    \item \textbf{Se houver correspondência:} substituir a palavra pela sua
    versão neutra correspondente. Preservar maiúsculas/minúsculas originais
    (ex.: se estava em maiúscula no início de frase, manter maiúscula).
    \item \textbf{Se não houver correspondência:} manter a palavra original.
    \item Recolocar a pontuação no lugar.
    \item Avançar para a próxima palavra.
  \end{enumerate}

  \item \textbf{Exibir o texto substituído:} Após processar todas as palavras,
  o texto resultante é apresentado com as substituições aplicadas.
\end{enumerate}
\end{algoritmo}

\bigskip

\subsection*{3.4 \ O Problema do ``Super-Homem'': Correspondência Parcial}

O enunciado faz uma pergunta muito perspicaz: \textit{``Como seu procedimento trataria
um termo como `super-homem'?''}

Este é um dos desafios clássicos de processamento de texto, chamado de
\textbf{problema de correspondência parcial} (\textit{substring matching}):
a palavra ``homem'' aparece \textit{dentro} de ``super-homem'', mas
``super-homem'' é uma palavra composta com significado próprio.

\textbf{Cenário 1 -- Algoritmo ingênuo (PROBLEMÁTICO):}

Se o algoritmo simplesmente procurar pela sequência de caracteres ``homem'' em
qualquer posição do texto e a substituir por ``pessoa'', o resultado seria:
\begin{center}
``super-homem'' $\to$ ``super-pessoa''
\end{center}
Isso pode ser aceitável em alguns contextos, mas perde o significado específico
da palavra composta.

\bigskip
\textbf{Cenário 2 -- Algoritmo aprimorado com correspondência por palavra completa:}

Um algoritmo mais sofisticado verificaria se a correspondência ocorre em uma
\textbf{palavra completa} (delimitada por espaços ou pontuação), não dentro de
outra palavra. Assim:
\begin{center}
Texto: ``O homem foi ao mercado'' $\to$ ``A pessoa foi ao mercado'' \\[4pt]
Texto: ``Ele é um super-homem'' $\to$ Reconhece ``super-homem'' como palavra
composta e consulta uma entrada específica no dicionário.
\end{center}

\textbf{Cenário 3 -- Dicionário com entradas compostas (MELHOR SOLUÇÃO):}

A melhor abordagem é incluir \textit{explicitamente} palavras compostas no dicionário
de substituições. A busca deve priorizar as entradas mais longas:

\begin{enumerate}
  \item Primeiro, verificar se a palavra/expressão atual corresponde a uma entrada
  composta no dicionário (ex.: ``super-homem'' $\to$ ``super-herói'').
  \item Somente se não houver correspondência composta, verificar palavras simples.
\end{enumerate}

\begin{atencao}
  O problema da correspondência parcial é um exemplo de por que a programação
  raramente é trivial. Soluções aparentemente simples frequentemente escondem
  casos especiais que precisam ser tratados. No Capítulo~8, você aprenderá as
  funções de manipulação de strings de C (\texttt{strstr}, \texttt{strcmp},
  \texttt{strcpy}), que tornam possível implementar esse algoritmo.
\end{atencao}

\bigskip

\subsection*{3.5 \ Algoritmo Completo Aprimorado}

\begin{algoritmo}
\textbf{Algoritmo: Substituição de Palavras com Gênero -- Versão Aprimorada}

\begin{enumerate}
  \item Carregar o dicionário de substituições na memória, organizando as entradas
  do \textit{maior} para o \textit{menor} número de palavras (para que palavras
  compostas sejam verificadas primeiro).

  \item Ler o parágrafo de texto completo do usuário.

  \item Percorrer o texto da esquerda para a direita, posição por posição.

  \item Para cada posição:
  \begin{enumerate}[label=\arabic*.]
    \item Tentar encontrar a \textbf{entrada mais longa} do dicionário que comece
    nessa posição.
    \item Verificar se a correspondência é com uma \textbf{palavra completa}
    (precedida e seguida por espaço, pontuação, início ou fim de texto).
    \item Se encontrar: realizar a substituição e avançar o número de posições
    correspondente ao tamanho da palavra encontrada.
    \item Se não encontrar: copiar o caractere atual para a saída e avançar
    uma posição.
  \end{enumerate}

  \item Exibir o texto resultante com todas as substituições aplicadas.
\end{enumerate}
\end{algoritmo}

\bigskip

\subsection*{3.6 \ Exemplo de Execução do Algoritmo}

\textbf{Entrada:}

\begin{tcolorbox}[colback=black!10, colframe=deitelblue, fontupper=\small]
``O homem e sua esposa foram ao mercado. Seu filho mais velho, que e um
verdadeiro super-homem, carregou todas as sacolas.''
\end{tcolorbox}

\textbf{Dicionário aplicado:}
``homem'' $\to$ ``pessoa''; \; ``esposa'' $\to$ ``cônjuge''; \;
``filho'' $\to$ ``filho(a)''; \; ``super-homem'' $\to$ ``super-herói''

\textbf{Saída esperada:}

\begin{tcolorbox}[colback=black!10, colframe=ecogreen, fontupper=\small]
``A pessoa e seu(sua) cônjuge foram ao mercado. Seu filho(a) mais velho, que e um
verdadeiro super-herói, carregou todas as sacolas.''
\end{tcolorbox}

\textbf{Observe} que ``super-homem'' foi tratado como uma única entrada do dicionário,
\textit{diferente} da entrada ``homem'', resolvendo exatamente o problema levantado
pelo enunciado.

\bigskip

\subsection*{3.7 \ Conexão com os Capítulos Futuros}

O enunciado menciona que você aprenderá no \textbf{Capítulo~4} que o termo formal
para ``procedimento'' é \textbf{algoritmo}. De fato:

\begin{tcolorbox}[colback=deitellightblue, colframe=deitelblue]
\textbf{Definição (antecipação do Capítulo~4):} Um \textbf{algoritmo} é uma sequência
finita e ordenada de passos bem-definidos que, quando executada, resolve um problema
específico em tempo finito.
\end{tcolorbox}

O processamento de substituições que descrevemos neste exercício é um algoritmo completo.
Ele especifica:
\begin{itemize}
  \item \textbf{O que fazer:} verificar cada posição do texto
  \item \textbf{A ordem:} da esquerda para a direita, entradas longas antes das curtas
  \item \textbf{A condição de parada:} ao atingir o fim do texto
\end{itemize}

Nos \textbf{Capítulos~3 e~4} você aprenderá a expressar algoritmos em pseudocódigo;
nos \textbf{Capítulos~7 e~8} você aprenderá ponteiros e manipulação de strings,
respectivamente -- as ferramentas de C necessárias para implementar este algoritmo
de forma eficiente.

\newpage

% ============================================================
\section*{Reflexão Final -- O Valor dos Exercícios ``Fazendo a Diferença''}
% ============================================================

\begin{tcolorbox}[colback=ecolightgreen, colframe=ecogreen,
  title={\bfseries Por que esses exercícios importam?}]

Os três exercícios deste capítulo -- a calculadora de emissão de carbono, a
calculadora de IMC e a neutralidade de gênero -- compartilham uma característica
essencial: são problemas \textbf{reais}, que afetam a vida real de pessoas reais.

\begin{itemize}
  \item \textbf{Exercício 1.10 (Carbono):} O aquecimento global é um dos maiores
  desafios da humanidade. Programas de computador são ferramentas poderosas para
  conscientizar e motivar mudanças de comportamento.

  \item \textbf{Exercício 1.11 (IMC):} A obesidade é uma epidemia global. Uma
  calculadora de IMC bem projetada pode ajudar milhões de pessoas a monitorarem
  sua saúde.

  \item \textbf{Exercício 1.12 (Gênero):} A linguagem molda a percepção social.
  Ferramentas de processamento de texto que promovem linguagem inclusiva contribuem
  para uma sociedade mais igualitária.
\end{itemize}

Esses problemas serão progressivamente implementados em C ao longo do livro.
O que você fez no Capítulo~1 foi \textbf{compreender o problema} -- a etapa mais
importante de qualquer projeto de software. Agora você está pronto para o
Capítulo~2, onde começará a escrever seus primeiros programas em C.
\end{tcolorbox}

\bigskip

\begin{boaprática}
  Um programador experiente compreende profundamente o problema que está resolvendo
  \textit{antes} de escrever uma única linha de código. Os exercícios ``Fazendo a
  Diferença'' treinam exatamente essa habilidade: a capacidade de analisar, modelar
  e planejar a solução -- independentemente da linguagem de programação usada.
\end{boaprática}

\end{document}
